\documentclass[12pt]{article}
\usepackage[margin=0.75in]{geometry}
\usepackage{fontspec}
\setmainfont{Latin Modern Roman}
\usepackage{microtype}
\usepackage{multicol}
\usepackage{fancyhdr}
\usepackage{titlesec}
\usepackage{enumitem}
\usepackage{xcolor}
\usepackage[hidelinks]{hyperref}
\setlength{\columnsep}{0.28in}
\setlength{\parindent}{1.1em}
\setlength{\parskip}{0.32em}
\setlength{\headheight}{15pt}
\titleformat{\section}{\large\bfseries\scshape}{}{0pt}{}
\titleformat{\subsection}{\normalsize\bfseries}{}{0pt}{}
\pagestyle{fancy}
\fancyhf{}
\fancyfoot[C]{\thepage}
\renewcommand{\headrulewidth}{0.4pt}
\newcommand{\focusbox}[1]{\par\vspace{0.4em}\noindent\begingroup\setlength{\fboxsep}{6pt}\colorbox{gray!12}{\begin{minipage}{\dimexpr\linewidth-2\fboxsep\relax}#1\end{minipage}}\endgroup\par\vspace{0.5em}}

\fancyhead[L]{Whately: Mood and Figure}
\fancyhead[R]{Student Reading}
\begin{document}
\begin{center}
{\LARGE\bfseries The Shape of an Argument}\par\vspace{0.25em}
{\large\scshape Week 5: Mood, Figure, and Formal Validity}\par\vspace{0.5em}
{Adapted and modernized from Richard Whately, \textit{Elements of Logic}}\par\vspace{0.6em}
\rule{0.82\textwidth}{0.4pt}
\end{center}

\section*{Central Question}
Can the shape of an argument tell us whether the conclusion follows?

\section*{Vocabulary Preview}
\begin{multicols}{2}
\begin{itemize}[leftmargin=*, itemsep=0.18em]
\item \textbf{Formal validity}: an argument's conclusion follows from its premises because of the argument's structure.
\item \textbf{Mood}: the A/E/I/O pattern of the major premise, minor premise, and conclusion.
\item \textbf{Figure}: the placement of the middle term in the two premises.
\item \textbf{A proposition}: universal affirmative: ``All S are P.''
\item \textbf{E proposition}: universal negative: ``No S are P.''
\item \textbf{I proposition}: particular affirmative: ``Some S are P.''
\item \textbf{O proposition}: particular negative: ``Some S are not P.''
\item \textbf{Distribution}: whether a term is being spoken of in its whole extent.
\item \textbf{Undistributed middle}: a faulty syllogism in which the middle term never covers the whole class it needs to connect.
\item \textbf{Illicit process}: a faulty syllogism in which a term is distributed in the conclusion but was not distributed in its premise.
\item \textbf{Affirmative conclusion}: a conclusion saying that one term belongs with another.
\item \textbf{Negative conclusion}: a conclusion saying that one term does not belong with another.
\end{itemize}
\end{multicols}

\section*{Introduction}
Week 4 taught you to build a practical syllogism map: major premise, minor premise, conclusion, major term, minor term, and middle term. That map showed whether the middle term seemed to connect the argument. Week 5 makes the map more exact.

Whately wants students of logic to notice that arguments have shape. Two arguments may be about completely different subjects and still have the same logical form. One may be about Socrates, one about political speeches, and one about Macbeth, yet all three may have the same pattern of universal and particular claims. If the form is valid, the conclusion follows whenever the premises are true. If the form is invalid, the conclusion may still be true by accident, but that particular argument has not proved it.

This week introduces mood and figure lightly. The goal is not to memorize every traditional valid form. The goal is to see how proposition type and term placement affect inference. A disciplined thinker should be able to say, ``This argument sounds forceful, but its form does not carry the conclusion.''

\focusbox{\textbf{Source note:} This reading is adapted and modernized from Whately's treatment of categorical syllogisms, especially his discussion of A/E/I/O propositions, mood, figure, and the rules that protect valid inference. The reading keeps Whately's formal concern while limiting the burden of memorization for a twelfth-grade long unit.}

\begin{multicols}{2}
\section*{Reading}

\subsection*{1. Why Form Matters}
A syllogism can be tested in two different ways. First, you may ask whether the premises are true. Second, you may ask whether the conclusion follows from those premises. These are not the same question.

Consider this argument:
\begin{itemize}[leftmargin=*]
\item All mammals are warm-blooded.
\item All whales are mammals.
\item Therefore, all whales are warm-blooded.
\end{itemize}

The argument is formally valid. Its shape carries the conclusion. If the premises are true, the conclusion must be true. Now replace the subject matter:
\begin{itemize}[leftmargin=*]
\item All unjust laws deserve resistance.
\item All tyrannical orders are unjust laws.
\item Therefore, all tyrannical orders deserve resistance.
\end{itemize}

The subject has changed, but the shape has not. Logic studies that shape.

\subsection*{2. Mood: The A/E/I/O Pattern}
Whately explains that the mood of a syllogism is found by naming the quantity and quality of its three propositions in order: major premise, minor premise, conclusion.

You already know the four basic proposition forms:
\begin{itemize}[leftmargin=*]
\item A: All S are P.
\item E: No S are P.
\item I: Some S are P.
\item O: Some S are not P.
\end{itemize}

In the whale example, the major premise is A, the minor premise is A, and the conclusion is A. Its mood is therefore A-A-A.

Mood matters because not every pattern can prove the conclusion it states. A syllogism with two negative premises, for example, does not connect the terms. If you say, ``No poets are careless thinkers'' and ``No historians are careless thinkers,'' you have not proved whether poets and historians are the same group, different groups, overlapping groups, or unrelated groups.

\subsection*{3. Figure: Where the Middle Term Stands}
Mood tells you what kind of propositions the argument uses. Figure tells you where the middle term stands in the premises.

Use this notation:
\begin{itemize}[leftmargin=*]
\item M = middle term
\item P = major term, the predicate of the conclusion
\item S = minor term, the subject of the conclusion
\end{itemize}

In the strongest beginner form, the middle term is subject of the major premise and predicate of the minor premise:
\begin{itemize}[leftmargin=*]
\item All M are P.
\item All S are M.
\item Therefore, all S are P.
\end{itemize}

Example:
\begin{itemize}[leftmargin=*]
\item All careful proofs state their reasons clearly.
\item All sound syllogisms are careful proofs.
\item Therefore, all sound syllogisms state their reasons clearly.
\end{itemize}

The middle term is ``careful proofs.'' It links the subject of the conclusion, ``sound syllogisms,'' to the predicate, ``things that state their reasons clearly.''

\subsection*{4. A Valid Form: A-A-A}
Traditional logicians named common valid forms so they could be remembered and tested. One famous form is often called \textit{Barbara}: A-A-A in the first figure.

Its structure is simple:
\begin{itemize}[leftmargin=*]
\item All M are P.
\item All S are M.
\item Therefore, all S are P.
\end{itemize}

This form is valid because whatever belongs inside M must also belong inside P, and S has been placed inside M. The conclusion only states the connection already forced by the premises.

You do not need to use the traditional name in ordinary writing. But you do need to recognize the principle: when the premises truly place S inside M and M inside P, S must be inside P.

\subsection*{5. An Invalid Form: The Undistributed Middle}
Now compare this tempting argument:
\begin{itemize}[leftmargin=*]
\item All careful readers notice contradictions.
\item All suspicious readers notice contradictions.
\item Therefore, all suspicious readers are careful readers.
\end{itemize}

This is invalid. The two groups may overlap, but the middle term, ``people who notice contradictions,'' has not been used widely enough to connect the two classes. The argument has said that careful readers belong among contradiction-noticers and that suspicious readers belong among contradiction-noticers. It has not proved that suspicious readers belong inside careful readers.

A picture may help. If two smaller circles both sit somewhere inside a larger circle, they are not necessarily the same smaller circle. They may overlap; they may not. The form has not forced the conclusion.

\subsection*{6. Negative and Particular Claims}
Mood also matters when claims become negative or particular.

If both premises are negative, the argument cannot connect the terms. ``No S are M'' and ``No P are M'' leaves S and P unconnected. Two separations from a third thing do not prove a relation between the first two things.

If one premise is negative, the conclusion must be negative. A negative premise separates a term from another term; it cannot produce an affirmative joining in the conclusion.

If one premise is particular, the conclusion usually must be particular. From a claim about ``some'' members of a class, you cannot usually prove a claim about ``all'' members of another class. Whately's rules protect the mind from taking a small amount of information and turning it into a larger conclusion than the premises permit.

\subsection*{7. The Form Test and the Fact Test}
Formal validity is powerful, but it is limited. Logic can say: ``If these premises are true, this conclusion follows.'' It cannot, by form alone, say that the premises are true.

That means a disciplined thinker asks two questions:
\begin{enumerate}[leftmargin=*]
\item \textbf{Form test:} Does the conclusion follow from the premises?
\item \textbf{Fact test:} Are the premises true, justified, or acceptable?
\end{enumerate}

Week 5 mostly trains the first question. Week 6 will turn more directly to the second question: how an argument can be valid in form and still fail because the premises are false, doubtful, or unproved.

\subsection*{8. Shakespearean Pressure}
Shakespeare's speakers often persuade by pressure rather than clean form. That does not mean they never reason. It means their reasoning must be patiently reconstructed.

Suppose someone argues: ``All honorable men love Rome. Brutus loves Rome. Therefore, Brutus is honorable.'' This sounds flattering, but the form is invalid. The first premise does not say that everyone who loves Rome is honorable. It says that honorable men love Rome. The middle term does not prove the conclusion.

Or consider Macbeth. Someone might argue: ``All fated events will happen. Macbeth's kingship will happen. Therefore, Macbeth's kingship is fated.'' Again, the argument is tempting but invalid. It treats ``will happen'' as if it automatically proves ``fated.'' The shape does not carry the conclusion.

\subsection*{9. What This Week Adds}
Week 4 asked whether an argument had a bridge. Week 5 asks whether the bridge has the right shape.

A mood-and-figure check should show:
\begin{itemize}[leftmargin=*]
\item the A/E/I/O form of each proposition;
\item the major, minor, and middle terms;
\item where the middle term appears in each premise;
\item whether the conclusion is too large, too positive, or too disconnected for the premises;
\item whether the argument passes a basic form test before the fact test begins.
\end{itemize}

This is not pedantry. Form protects judgment. It keeps readers from accepting a conclusion merely because the subject matter is familiar, the speaker is confident, or the conclusion already feels attractive.

\section*{Reading Focus}
\begin{enumerate}[leftmargin=*]
\item What is the difference between testing an argument's form and testing its facts?
\item Define mood in your own words, and give one example using A/E/I/O letters.
\item Define figure in your own words. What does figure track that mood does not track?
\item Why is A-A-A in the first figure a valid beginner form?
\item Explain the flaw in this form: All P are M. All S are M. Therefore, all S are P.
\item Why can two negative premises fail to prove a conclusion?
\item Choose either the Brutus or Macbeth example. What makes the argument sound persuasive, and what makes its form questionable?
\end{enumerate}

\section*{Handout Summary}
Write 5--7 sentences explaining how mood and figure help a reader test formal validity. Your summary must include the words \textit{mood}, \textit{figure}, \textit{middle term}, and \textit{follows}.
\end{multicols}
\end{document}
